\chapter{Bitstoppning och CRC-15}
\label{ch:stuffing}

{\large\itshape De två mekanismerna som gör att en frame överlever en störig buss.\par}

\begin{chapterbox}{I det här kapitlet}
\begin{itemize}
  \item Varför en lång rad lika bitar är ett problem, och femregeln som löser det.
  \item Vilka fält som stoppas, och vilka som inte gör det.
  \item CRC-15: polynomet, den bitseriella motorn, och ett exempel räknat för hand.
  \item Varför samma motor både beräknar och kontrollerar.
\end{itemize}
\end{chapterbox}

\section{Problemet med lika bitar}

CAN-noder har ingen gemensam klocka. Varje nod har sin egen oscillator, och de går aldrig exakt
lika fort. En mottagare håller takten genom att titta på \emph{flankerna} i bitströmmen: varje
gång signalen byter nivå vet den var en bitgräns låg, och kan justera sin egen räknare
(\kapref{ch:tajming}).

Kommer det inga flanker driver mottagaren i väg. Efter tillräckligt många lika bitar i rad
samplar den fel bit, och framen är förlorad.

\textbf{Lösningen är att göra flankerna obligatoriska.} Efter fem lika bitar i rad sätter
sändaren in en \textbf{stoppbit} med motsatt värde, som mottagaren känner igen och kastar.

\section{Femregeln}

Regeln är kort: \textbf{efter fem lika bitar i rad sänds en bit med motsatt värde}, oavsett vad
som kommer härnäst.

Ett exempel, med den insatta biten markerad:

\begin{console}
in:  1 1 1 1 1 0 0 0 0 0 1
ut:  1 1 1 1 1 0 0 0 0 0 1 0 1
               ^           ^
           stoppbit    stoppbit
\end{console}

Det finns två detaljer här som är lätta att få fel, och båda är klassiska buggar:

\textbf{Stoppbiten räknas själv.} Efter att en stoppbit satts in börjar räkningen om med
stoppbiten som första bit i den nya serien. Den är en helt vanlig bit i strömmen, för alla
utom den som ska plocka bort den igen.

\textbf{Regeln gäller strömmen, inte fältet.} Fem lika bitar som sträcker sig över gränsen
mellan identifieraren och DLC:n utlöser en stoppbit precis som fem inom ett fält. En
implementation som nollställer räknaren vid varje fältbyte fungerar nästan alltid, och går
sönder precis när två fält råkar sluta och börja med samma bitvärde.

\subsection{Vilka fält som stoppas}

Stoppning gäller från SOF till och med CRC-fältet. Därefter \textbf{slutar} den:

\begin{narrowtable}{@{}ll@{}}
\thead{Fält} & \thead{Stoppas?} \\ \midrule
SOF, arbitrering, kontroll, data, CRC & Ja \\
CRC-avgränsare, ACK, EOF, intermission & Nej \\
\end{narrowtable}

Att slutet inte stoppas är vad som gör EOF användbart som slutmärke. Sju recessiva bitar i rad
kan inte förekomma i den stoppade delen --- där är fem lika det mesta som får passera --- så en
mottagare som ser sju recessiva vet säkert att framen är slut.

\subsection{Stoppfel}

Mottagaren räknar likadant som sändaren. Efter fem lika bitar \emph{måste} nästa bit vara den
motsatta. Är den inte det är något fel, och mottagaren rapporterar ett \textbf{stoppfel}.

Det är ett av CAN:s billigaste feldetekteringsmedel: det kostar ingen extra information alls,
bara en räknare i vardera änden.

\begin{aside}
\note[Kostnaden] Stoppningen gör framen längre, och hur mycket beror på datan. En frame vars
databitar växlar hela tiden får inga stoppbitar alls; en frame full av nollor får en var femte
bit. För den längsta dataframen är påslaget som mest ett tjugotal bitar. Det är också därför en
CAN-frames längd inte är exakt förutsägbar från DLC:n --- vilket i sin tur är därför den som
räknar på bussbelastning räknar med värstafallet.
\end{aside}

\section{CRC-15}

Checksumman är femton bitar, beräknad över allt från SOF till och med datafältet --- alltså de
\emph{verkliga} bitarna, utan stoppbitar. Stoppbitar är transportens sak och ingår inte i det
som skyddas.

\subsection{Polynomet}

CAN använder

\[ x^{15} + x^{14} + x^{10} + x^{8} + x^{7} + x^{4} + x^{3} + 1 \]

som i binär form, koefficienterna för $x^{14}$ ned till $x^{0}$, blir

\begin{console}
100010110011001
\end{console}

Polynomet är valt för att upptäcka alla enkel-, dubbel- och trippelfel, alla felburstar upp till
15 bitar, och de allra flesta längre. Det är inte en slumpmässig konstant.

\subsection{Motorn}

Beräkningen är ett femtonbitars skiftregister med återkoppling, och den tar en bit i taget:

\begin{cppcode}
// One bit into the CRC-15 engine. crc is 15 bits wide.
const bool feedback{static_cast<bool>((crc >> 14U) & 1U) ^ bit};
crc = static_cast<std::uint16_t>((crc << 1U) & 0x7FFFU);
if (feedback)
{
    crc ^= 0x4599U;   // 100010110011001
}
\end{cppcode}

Tre rader: XOR:a den inkommande biten med registrets översta bit, skifta registret ett steg, och
XOR:a in polynomet om resultatet av den första operationen var ett.

\subsection{Ett exempel för hand}

Sekvensen \code{10110010}, matad in bit för bit i ett nollställt register:

\begin{narrowtable}{@{}llll@{}}
\thead{Bit} & \thead{Värde} & \thead{Registret efter} & \thead{Hex} \\ \midrule
1 & 1 & \code{100010110011001} & \code{0x4599} \\
2 & 0 & \code{100111010101011} & \code{0x4EAB} \\
3 & 1 & \code{001110101010110} & \code{0x1D56} \\
4 & 1 & \code{111111100110101} & \code{0x7F35} \\
5 & 0 & \code{011101111110011} & \code{0x3BF3} \\
6 & 0 & \code{111011111100110} & \code{0x77E6} \\
7 & 1 & \code{110111111001100} & \code{0x6FCC} \\
8 & 0 & \code{001101000000001} & \code{0x1A01} \\
\end{narrowtable}

Checksumman är alltså \code{0x1A01}.

\subsection{Samma motor kontrollerar}

Här ligger elegansen. Mottagaren har ingen separat kontrollfunktion: den matar in \emph{samma}
bitar plus \emph{den mottagna checksumman} genom samma motor.

Om allt stämmer står registret på \textbf{noll} när sista CRC-biten passerat. Matar vi in
\code{10110010} följt av \code{001101000000001} ur exemplet ovan blir slutvärdet
\code{0x0000}.

Det är en egenskap hos konstruktionen, inte ett specialfall: att lägga resten sist i dividenden
gör den jämnt delbar. Sändare och mottagare kan därför använda exakt samma modul, utan
lägesomkoppling --- vilket är precis vad en hårdvarukonstruktion vill ha.

\begin{aside}
\note[En detalj som bits] Motorn behöver ändå kunna nollställas. En avbruten frame lämnar skräp i
registret, och nästa frame som börjar med ett smutsigt register får fel checksumma från första
biten. Nollställningen sker vid varje framestart.
\end{aside}

\subsection{Vad CRC:n inte skyddar}

CRC:n täcker SOF till och med datafältet. Den täcker alltså \emph{inte} CRC-avgränsaren,
ACK-fältet eller EOF. Fel i dem upptäcks i stället av formatreglerna: en avgränsare som inte är
recessiv är ett formatfel (\kapref{ch:fel}).

\section*{Övningar}

\exercise{Stoppa en ström}{Räkning}
\label{ex:3:1}
Stoppa följande bitströmmar enligt femregeln, och markera de insatta bitarna.
\begin{enumerate}[a)]
  \item \code{0000011111}
  \item \code{1111111111}
  \item \code{0101010101}
\end{enumerate}
Hur många bitar blev var och en efter stoppning?

\exercise{Avstoppa}{Räkning}
\label{ex:3:2}
En mottagare tar emot \code{111110000010}.
\begin{enumerate}[a)]
  \item Vilka bitar är stoppbitar?
  \item Vilken var den ursprungliga bitströmmen?
  \item Antag i stället att mottagaren fått \code{111111...}. Vad rapporterar den, och efter
    vilken bit?
\end{enumerate}

\exercise{CRC-motorn}{Räkning}
\label{ex:3:3}
Mata in \code{1101} i ett nollställt CRC-15-register, ett steg i taget, och skriv upp registrets
värde efter varje bit. Använd tabellen i avsnitt 3.3 som mall.

\exercise{Varför noll?}{Resonemang}
\label{ex:3:4}
\begin{enumerate}[a)]
  \item Varför står mottagarens register på noll när både data och checksumma matats in?
  \item Vilken praktisk fördel har det för en hårdvarukonstruktion?
  \item En nod matar av misstag in \emph{stoppbitarna} i CRC-motorn också. Vad händer, och vem
    upptäcker det?
\end{enumerate}
